\subsection{Contrast-specific minimum detectable effects}\label{sec:mde}
The regional score scales differ by contrast, not only by region. At 90-day blocks the two-sided 95 per cent minimum detectable effect is 9.24~nats for New Zealand shared attenuation with free trend, but only 5.42~nats for the zero-trend arm and 2.74~nats for the fixed-null-trend arm (against $-17.62$, $-11.43$ and $-5.01$~nats respectively). The fixed-null loss is $-3.59$ of its own block standard error and $1.83$ times its own 90-day MDE, not below the free-trend MDE against which it was previously compared. The other contrasts are 3.82~nats for California \texttt{LI+} total attenuation (against $-2.39$~nats observed; the pooled percentile half-width is 3.27~nats), 48.30~nats for California \texttt{EP+} minus \texttt{EP} (against $+69.55$~nats at the binding bound and $+77.06$~nats at the interior scale), 43.35~nats for California \texttt{OC} minus \texttt{EP} (against $-51.05$~nats), and 47.94~nats for the three-year Japan \texttt{OC} minus \texttt{EP} total (15.98~nats per three-year fold, against $-221.21$~nats total and $-73.74$~nats per fold). The nine genuine annual JMA blocks, pooled within the three-year folds, give $\mathrm{SE}_{365}=28.83$~nats and $\mathrm{MDE}_{95}=56.51$~nats. The California \texttt{LI+} total deficit is below its own 90-day detection scale, whereas the New Zealand fixed-null deficit and the Japan total exceed theirs. The completed \texttt{EP+} gain is of the same order as its own contrast-specific MDE, not the much larger attenuation MDE. These calculations delimit what the negative and positive summaries can detect; they do not convert a non-significant contrast into evidence of physical absence. Table~\ref{tab:mde} expands every regional contrast at all three block lengths, and the \texttt{supplement} records the dependent-block qualification.

\begin{table}[htbp]\centering\scriptsize
\setlength{\tabcolsep}{2pt}
\resizebox{\textwidth}{!}{%
\begin{tabular}{@{}lrlrrrr@{}}\toprule
Contrast & $N_{\rm test}$ & Observed $\Delta\ell$ (nats) & Block (d) & $\mathrm{SE}_B$ (nats) & $\mathrm{MDE}_{95}$ (nats) & $\mathrm{MDE}_{95}$/event\\\midrule
NZ shared $-$ ML & 612 & $-17.623$ free $d$; $-11.431$ zero $d$; $-5.006$ fixed null $d$ & 90 & 4.71 & 9.24 & 0.0151 \\
 & & & 180 & 5.42 & 10.62 & 0.0174 \\
 & & & 365 & 6.65 & 13.04 & 0.0213 \\
\addlinespace
CA \texttt{LI+} total $-$ floor & 315 & $-2.387$ & 90 & 1.95 & 3.82 & 0.0121 \\
 & & & 180 & 1.94 & 3.81 & 0.0121 \\
 & & & 365 & 1.94 & 3.81 & 0.0121 \\
\addlinespace
CA \texttt{EP+} $-$ \texttt{EP} & 315 & $+69.551$ bound; $+77.056$ interior $\sigma_0$ & 90 & 24.64 & 48.30 & 0.1533 \\
 & & & 180 & 24.17 & 47.38 & 0.1504 \\
 & & & 365 & 26.96 & 52.84 & 0.1678 \\
\addlinespace
CA \texttt{OC} $-$ \texttt{EP} & 315 & $-51.051$ & 90 & 22.12 & 43.35 & 0.1376 \\
 & & & 180 & 21.97 & 43.06 & 0.1367 \\
 & & & 365 & 21.50 & 42.15 & 0.1338 \\
\addlinespace
Japan \texttt{OC} $-$ \texttt{EP} (per 3-year fold) & 1277 & $-73.737$ & 3-year & 8.15 & 15.98 & 0.0375 \\
Japan \texttt{OC} $-$ \texttt{EP} (3-fold sum) & 1277 & $-221.212$ & 3-year & 24.46 & 47.94 & 0.0375 \\
Japan \texttt{OC} $-$ \texttt{EP} (9 annual blocks) & 1277 & $-221.212$ & 365 & 28.83 & 56.51 & 0.0442 \\
\bottomrule\end{tabular}}
\caption{Contrast-specific minimum detectable effects at two-sided ninety-five per cent. Each row names the exact scored contrast, its test-event count, the observed pooled total and the block length. $\mathrm{SE}_B$ is the independent-block standard error of paired calendar-block score totals; $\mathrm{MDE}_{95}=1.96\,\mathrm{SE}_B$ and per-event values divide by $N_{\rm test}$. New Zealand and California use 90-, 180- and 365-day blocks. The Japanese archive contains three three-year folds (2016--2018, 2019--2021, 2022--2024), not three annual fold totals. Its per-fold row reports the between-window standard error 8.15~nats, the three-fold-sum row reports $3\times8.15=24.46$~nats, and the final row pools the nine genuine annual JMA blocks within folds, giving $\mathrm{SE}_{365}=28.83$~nats and $\mathrm{MDE}_{95}=56.51$~nats. Contiguous dependent blocks are descriptive dependence-scale units, not independent replications.}
\label{tab:mde}
\end{table}

The corresponding fractional rate perturbations $\delta=\Delta\lambda/\lambda$ from the quadratic small-perturbation score approximation are 17.4, 18.6 and 20.6\% for the New Zealand free-trend arm, 13.3, 13.4 and 13.9\% for the zero-trend arm, 9.5, 9.6 and 10.0\% for the fixed-null arm, 55.4, 54.8 and 57.9\% for California \texttt{EP+}, and 27.4\% for the per-fold and three-fold Japan summaries (29.7\% for the nine annual blocks). The 0.0077~nat-per-event New Zealand scale is much finer than the 0.0782~nat-per-event California \texttt{EP+} scale because California's block scores are more dispersed. The completed \texttt{EP+} gain of 69.551~nats is of the same order as its own California 90-day MDE, while the S8 raw mean of 0.294~nats is far below its own MCSE.
